\chapter{Background Research}
\label{ch:background}

% Rebuilt 2026-08-02 from the frozen sources per DISSERTATION_ARCHITECTURE.md §Ch2 and
% AUDIT_DOSSIER.md Part 4. Opening adapted from the June ch2 opening (quarry-safe,
% dossier Part 3); all section blocks assembled from the W1 drafts, seams smoothed.
% Bib fragments merged into refs.bib at assembly (commit a923b53); NEWREF entries pending W3 verification.

The two explanations Chapter~\ref{ch:intro} left entangled have been examined before, in pieces, by five research strands that rarely engage on equal terms; each is read here for the evaluation standard it brings to positive results for text, because that standard, as much as the data or the model, decides the verdict. The econometrics of volatility forecasting supplies the price-based models that set the reference any textual claim must improve upon; the text-in-finance literature makes the claim under test. Neural representation and the long-document problem supply the machinery for testing it, and prompted large language models are its newest form. Shortcut learning and evaluation hygiene, imported from outside finance, give the investigation its audit character. The closing section turns that reading into the choices this investigation adopts.

\section{Literature Survey}
\label{sec:bg-survey}

Each strand closes by recording what is established, what is contested, and the gap this project addresses.

\subsection{Price-Based Volatility Forecasting}
\label{sec:bg-price}

Volatility, as used throughout this report, is the conditional standard deviation of an asset's returns. Returns themselves are notoriously close to unpredictable; volatility is not. Calm days cluster with calm days and turbulent with turbulent, and that persistence makes volatility one of the few genuinely forecastable quantities in financial markets. Engle's autoregressive conditional heteroskedasticity (ARCH) model~\cite{engle1982arch} first formalised the regularity, volatility clustering, by letting today's return variance depend on the squared magnitudes of recent shocks. Bollerslev~\cite{bollerslev1986garch} generalised the recursion by adding lagged variances, yielding the GARCH family, whose simplest member, GARCH(1,1), captures volatility persistence with only three parameters % src: bollerslev1986garch (literature)
and remains the parametric workhorse four decades later. Nelson's EGARCH~\cite{nelson1991egarch} models the log variance and permits an asymmetric response, capturing the leverage effect: negative returns raising future volatility more than equal positive ones.

Risk managers, option desks and position-sizers need a volatility forecast for the coming days and weeks, not for the coming years. Hence the forecast horizons of 5, 10 and 20 trading days used throughout the project, % src: frozen writing/dissertation/supporting/sections/03_dataset.tex (h in {5,10,20})
and hence the evaluation standard. A forecast at these horizons is useful only if it holds strictly out of sample, under the temporal ordering a deployed system would face.

A second paradigm treats volatility as approximately observable rather than latent. Andersen et al.~\cite{andersen2003modeling} showed that realised volatility (RV), built by aggregating ever-finer intraday squared returns, estimates a day's variation consistently, and that simple regressions on realised measures can match or beat latent-variance GARCH out of sample. Forecasting thereby becomes an ordinary regression on an observed target, and this project adopts the reframing, which lets price and text predictors compete within one supervised-learning protocol. Corsi's heterogeneous autoregressive model, HAR-RV~\cite{corsi2009har}, became the workhorse of that tradition by regressing future RV on its daily, weekly and monthly averages. Three slope coefficients, % src: corsi2009har (literature)
estimable by ordinary least squares, reproduce the long-memory signature of volatility, and they are exceptionally difficult to beat.

That difficulty is among the best-replicated findings in the field. Comparing 330 ARCH-type specifications, % src: hansen2005garch (literature)
Hansen and Lunde~\cite{hansen2005garch} found no evidence that anything outperformed a plain GARCH(1,1) for exchange-rate volatility. The survey of Poon and Granger~\cite{poon2003forecasting}, covering 93 published studies, % src: poon2003forecasting (literature)
was similarly sobering. Simple specifications forecast about as well as complex ones, and apparent improvements often dissolve out of sample. Nor does outside information straightforwardly help. Paye~\cite{paye2012dejavol} found that macroeconomic predictors, however correlated in sample, rarely deliver significant out-of-sample gains over autoregressive benchmarks. A new information source therefore has to do more than predict volatility in isolation; it has to add something \emph{beyond} a strong price-based benchmark. Blair et al.~\cite{blair2001forecasting} exemplify the incremental framing, asking how much implied volatility and intraday returns each contribute once the other sources are conditioned upon.

Plain HAR-RV is not the end of the price-only frontier. Patton and Sheppard~\cite{pattonsheppard2015shar} split realised variance into signed semivariances (negative and positive returns taken separately) and find the downside component carries most of the predictive content, so semivariance-augmented specifications typically improve on the plain model. Bollerslev et al.~\cite{bollerslev2016harq} discount days whose realised measure is noisily estimated, and a long tradition adds option-implied volatility (the market's own assessment, as in the VIX) as an exogenous regressor~\cite{blair2001forecasting}. Neural price models are reported to beat the HAR lineage~\cite{bucci2020neural,christensen2023mlvol}, and none was fitted here (Section~\ref{app:external:neural}). ``The'' price benchmark is therefore a moving target. A text model beating one specification may still lose to a stronger neighbour, so a credible audit must allow the price side its best representative. In synthesis: volatility is forecastable; HAR-RV and GARCH(1,1) are strong, parsimonious references; exogenous information rarely adds out-of-sample value; and the question inherited here is whether regulatory-disclosure text, held to that standard, is an exception.

\subsection{Text and Disclosures in Finance}
\label{sec:bg-text}

The claim that language carries market-relevant information predates modern language technology. Tetlock~\cite{tetlock2007giving} quantified the pessimism of a daily \emph{Wall Street Journal} column by counting words and found that high pessimism predicts downward pressure on prices followed by reversion. Message-board sentiment tracks indices while forecasting poorly stock by stock~\cite{das2007yahoo}, and a century of news columns % src: garcia2013sentiment (literature)
shows the predictive content of sentiment concentrating in recessions~\cite{garcia2013sentiment}: evidence that text signals aggregate better than they forecast and are regime-dependent.

A parallel methodological line asks how financial text should be measured. Loughran and McDonald~\cite{loughran2011liability} demonstrated that general-purpose sentiment dictionaries are badly miscalibrated for financial documents: almost three-quarters of the words counted as negative by the widely used Harvard psychosocial dictionary (\emph{liability}, \emph{tax}, \emph{cost}) are not negative in a 10-K. % src: loughran2011liability (literature)
They also constructed the finance-specific word lists that remain the standard resource. At the aggregate level, news text demonstrably carries economic content, proxying risk perceptions and tracking the business cycle~\cite{bybee2024business,manela2017nvix}. Yet Loughran and McDonald's survey~\cite{loughran2016survey} issues warnings this project treats as design requirements: textual measures are noisy, document length and boilerplate confound naive counts, and machine-learning approaches invite overfitting whenever they are judged in sample or against weak baselines.

Regulatory disclosures are a distinct and more promising class of text than journalist-mediated news. United States public companies must file an annual report (Form 10-K) and quarterly reports (Form 10-Q): book-length documents containing financial statements, management's discussion of results and enumerated risk factors. A current report (Form 8-K) is due within days of any material event. Filings are primary sources written by the firm under legal liability, timestamped precisely on the Securities and Exchange Commission's EDGAR system, and available for the full cross-section of listed firms. The same legal regime shapes the language: litigation risk encourages hedged, exhaustive, boilerplate prose, and Cohen et al.~\cite{cohen2020lazy} document that the bulk of a periodic filing repeats its predecessor nearly verbatim, with the \emph{changes} carrying the predictive content. Whatever signal a periodic filing holds is spread thinly through tens of thousands of words.

The founding study of volatility prediction from filings is Kogan et al.~\cite{kogan2009predicting}, who regress firm return volatility on word features of the 10-K management-discussion section, using support vector regression over TF--IDF word weights (a scheme that discounts ubiquitous words). Text combined with a historical-volatility feature outperformed that feature alone, and the result is widely cited as evidence that disclosures predict risk. Its evaluation design, however, merits scrutiny: the baseline is a single lagged-volatility regressor rather than a HAR-class model, and the horizon is annual, not the daily-to-monthly range risk applications require. Significance against the baseline rests on a permutation test that treats filings as independent observations, the observation-level convention whose day-clustered repair Chapter~\ref{ch:data} adopts. The empirical chapters return to this design directly, reproducing its evaluation recipe on the present benchmark and then repairing its elements one at a time. % forward pointer to the replicate-and-dissolve programme (results/tables/kogan_corpus_audit.md, kogan_dissolve.md); no numbers stated here

Later work upgrades the representations and, unevenly, the evaluation. FinText~\cite{rahimikia2021fintext} trains word embeddings (dense vector representations of words learned from co-occurrence patterns) on financial newswire text and applies them to realised-volatility forecasting, benchmarking against HAR-family models. It is the standing counterexample to any blanket claim that textual studies ignore strong price baselines. FETILDA~\cite{fetilda2022} fine-tunes long-document neural encoders on 10-K volatility regression, bringing modern architectures to bear on precisely the Kogan task. A neighbouring literature studies earnings calls, the calls at which executives present results and take analyst questions. MDRM~\cite{qin2019mdrm} predicts post-call volatility from verbal and vocal cues jointly, and later corpus work scaled the setting~\cite{maec2020}.

Most consequential for the present investigation, Yu et al.~\cite{yu2025samecompany} establish that transcript embeddings from the same firm are markedly more similar than embeddings from different firms, and that an identity-only baseline (a model told nothing except which firm is speaking) already matches or beats transcript-based models. This is the direct precursor of the identity concern at the centre of this project. Because firms file repeatedly and volatility is strongly firm-persistent, an expressive representation can succeed by encoding \emph{which firm is speaking} rather than \emph{what is said}. That is an instance of the shortcut learning surveyed in the final strand~\cite{geirhos2020shortcut}. What Yu et al. leave open is the forecast-level verdict: whether, under a calibrated protocol with a strong recalibrated reference, text still adds value once firm identity is controlled. Building the instrument that can return that verdict is what this project attempts.

In synthesis: financial text demonstrably correlates with prices and risk, domain-appropriate measurement matters, and one modern design (FinText) benchmarks against the HAR family. What no surveyed design imposes is the conjunction adopted here: a \emph{recalibrated} price reference, one whose biases are corrected on held-out data before any comparison; a firm-identity control, asking whether text beats a reference that knows only who is filing; and significance tests that treat the trading day, not the filing, as the independent unit. The gap is one of protocol rather than of baselines, and Section~\ref{app:external:benchmarks} tabulates it corpus by corpus. % characterisation matches frozen writing/dissertation/supporting/sections/02_related.tex

\subsection{Neural Representations and the Long-Document Problem}
\label{sec:bg-neural}

Dictionaries and bag-of-words features treat words as independent counters, blind to context, negation and syntax. The Transformer of Vaswani et al.~\cite{vaswani2017attention} addressed this with self-attention, in which every token's representation is updated as a weighted combination of all others, capturing long-range dependencies at a computational cost quadratic in sequence length. Devlin et al.~\cite{devlin2019bert} showed that a bidirectional Transformer encoder (BERT), pretrained on large unlabelled corpora and then \emph{fine-tuned} (its weights further trained on the downstream task), transfers across tasks with minimal task-specific machinery. BERT fixed a maximum input of 512 tokens, % src: devlin2019bert (literature)
a constraint inherited by most descendants. RoBERTa~\cite{liu2019roberta}, identical in structure but pretrained longer on more data, outperformed BERT across benchmarks: a caution against attributing performance differences between encoders to architecture alone.

Fine-tuning adapts on the order of a hundred million parameters % src: devlin2019bert (BERT-base size, literature)
to a noisy scalar target, an extreme capacity-to-signal ratio that invites the encoder to memorise idiosyncrasies of the training period. Freezing the encoder and training only a small head trades that adaptation for stability. On that trade-off for financial regression the literature offers little guidance, its benchmarks dominated by classification. Domain adaptation follows the same pattern: the two models sharing the name FinBERT (Araci's~\cite{araci2019finbert}, further pretrained on financial news, and Yang et al.'s~\cite{yang2020finbert}, pretrained from scratch on financial communications) demonstrate clear gains on short-text sentiment \emph{classification}. Evidence that domain pretraining helps on \emph{regression} over book-length documents against strong numerical baselines is far scarcer.

The 512-token window collides with the documents this project studies. The long-form filings of the benchmark assembled in Chapter~\ref{ch:data} have a median length of 29{,}278 wordpiece tokens % src: frozen 03_dataset.tex Table 1 (median long-form length)
(sub-word units), and every one exceeds 4{,}096 tokens. % src: frozen 03_dataset.tex Table 1 (share >4,096 tokens = 100%)
A 512-token encoder therefore reads roughly two per cent of the median filing, and a 4{,}096-token window roughly fourteen per cent. % src: TECHNICAL_SUPPLEMENT.md token arithmetic (S1 ~1.75--2%, S5 ~14%)
The literature offers four families of response. \emph{Truncation} keeps only the first window of text: cheap, but a 10-K's mandated structure gives no reason to expect the predictive content in its opening pages. \emph{Chunking} splits the document into window-sized segments, encodes each independently and pools the embeddings, by unweighted averaging or by learned attention (close to sixty would be needed to read a median filing exhaustively; the arms executed here cap at eight). % src: TECHNICAL_SUPPLEMENT.md (~57--60 chunks per median filing); configs/models/C1_bert_s2.yaml, C2_finbert_s\{2,3,4\}.yaml (max\_chunks: 8, chunk\_stride: 256)
\emph{Hierarchical} models make the aggregation itself a learned sequence model over chunk embeddings~\cite{pappagari2019hierarchical}. \emph{Sparse attention} redesigns the attention pattern: Longformer~\cite{beltagy2020longformer} combines a sliding window with a few global tokens, cutting complexity from quadratic to linear and enabling a native context of 4{,}096 tokens, % src: beltagy2020longformer (literature)
while remaining initialisable from RoBERTa weights. This isolates the effect of context length from that of pretraining lineage. These responses, with the two pooling variants distinguished, are fixed as strategies S1--S5 in the design of Chapter~\ref{ch:data}.

Since 2020 the representation frontier has moved beyond encoders of a few hundred million parameters. Decoder-only large language models (generative models of billions of parameters, such as the open-weight Qwen3 family~\cite{qwen3_2025}) can be repurposed as \emph{frozen} text embedders, their hidden states pooled into a document vector with no task-specific fine-tuning~\cite{wang2024llmembed}. The regime is representationally richer than the encoders above yet immune by construction to fine-tuning overfitting, and the project's model spectrum includes three such embedders of seven to eight billion parameters. % src: frozen 04_methods.tex (C5 = three frozen 7--8B embedders)
Prompting the same models to produce forecasts directly raises distinct evaluation questions, surveyed in the next strand.

Read critically, this strand carries two caveats. First, even the longest standard contexts fall sevenfold short of the median long-form filing: ``long-context'' models still read a minority of each filing. The operative question is therefore whether reading eight times more text % src: 4,096/512 (literature arithmetic; beltagy2020longformer, devlin2019bert)
helps, not whether the document is read in full. Second, published comparisons of these strategies sit almost exclusively on classification benchmarks. FETILDA~\cite{fetilda2022} is the noteworthy exception on financial regression, but no existing comparison, FETILDA's included, is run against a recalibrated strong price reference with formal significance testing and firm-identity controls. In synthesis: pretrained Transformer encoders are the strongest general-purpose text representations; several viable long-document strategies exist, compared almost exclusively on classification benchmarks; and frozen billion-parameter embedders extend the representational spectrum upwards. What the literature leaves unsettled, whether any point on that spectrum adds forecast value over the strongest price references under an audit-grade protocol, is the question Chapters~\ref{ch:results} and~\ref{ch:validation} answer.

\subsection{Prompted Language Models as Forecasters}
\label{sec:bg-prompted}

Every strand surveyed so far assumes a text model must be \emph{trained} on the forecasting task before it can forecast. A strand that emerged once conversational large language models became publicly available in late 2022 removes that assumption. An instruction-tuned large language model (LLM) is a neural network pretrained on web-scale text and then further adjusted to follow written instructions. \emph{Prompting} means handing such a model a document and a natural-language request (in effect, ``read this filing and state how volatile the firm's shares will be'') and reading a forecast out of its textual reply, with no task-specific training at all. The strand's reference result is due to Lopez-Lira and Tang~\cite{lopezlira2023chatgpt}, who prompted ChatGPT to judge whether news headlines were good or bad news for the firm concerned. The resulting scores predicted next-day stock returns while earlier and smaller models given the same instruction produced no comparable signal, making prompting a serious candidate forecasting technology.

A prompted forecast nevertheless has a different epistemic status from a trained one, and the evaluation literature has identified three threats. The first is \emph{data contamination}: the material a model is evaluated on may already sit inside its pretraining corpus. Sainz et al.~\cite{sainz2023contamination} document contamination across standard benchmarks and argue that it must be measured, benchmark by benchmark, rather than assumed away. In a forecasting setting the threat is sharper, because contamination operates as look-ahead: a web-scale corpus contains text written \emph{after} the events to be forecast, so a model asked to ``predict'' volatility for a date inside its pretraining window may simply be recalling what happened. Two defences follow. An \emph{era split} compares performance before and after the model's training-data cutoff, the date at which its corpus ends, since outcomes beyond it cannot have been memorised. \emph{Zero-content probes} strip the document from the prompt, retaining only the metadata header (form type, item codes, filing date) or that header plus the stock ticker, so that any skill which survives measures recall and firm recognition rather than reading.

The second threat is \emph{elicitation fragility}. What a prompted model returns depends on the exact wording and format of the request and on decoding settings such as temperature, which governs how much randomness enters word-by-word generation. A gain that disappears when the instruction is paraphrased is a property of one prompt rather than of the model. A nominally deterministic temperature-zero configuration must still be checked for repeat-decode stability, and a cross-family comparison is fair only if each family receives its own prompt adaptation before being declared unable to do the task. The third threat is \emph{instrument degeneracy}: a prompted model can stop acting as a forecaster, most visibly through mode collapse, returning one value for most inputs whatever it reads. Screening such runs out is legitimate only if the screen reads the forecast's own dispersion and accuracy, never the contrast under test, else screening becomes selection on the hoped-for conclusion.

For the disclosure--volatility question these considerations expose a specific gap. No study in the lineage of Section~\ref{sec:bg-text} has placed a prompted forecaster inside the same benchmark, chronological splits and reference forecasts as trained representations of the identical text. Whether prompting and fine-tuned encoding extract the same information is therefore unknown. The design of Chapter~\ref{ch:data} evaluates a prompted arm under the protocol every trained model faces, and Chapter~\ref{ch:validation} applies the battery this strand demands: contamination probes, era splits, paraphrase and repeat-decode checks, and cross-family replication.

\subsection{Shortcut Learning and Evaluation Hygiene}
\label{sec:bg-shortcut}

The final strand is not financial at all, yet it supplies the intellectual frame for this project's central threat. Geirhos et al.~\cite{geirhos2020shortcut} call \emph{shortcut learning} the failure mode in which a network scores well by exploiting an incidental feature correlated with the label rather than the competence the benchmark meant to measure. Ordinary held-out evaluation therefore cannot tell the two apart. The most instructive instances concern entities. In clinical machine learning, models evaluated with record-wise splits (records from one patient falling on both sides of the train--test divide) learn to recognise the patient rather than the disease, and accuracy collapses once splitting is subject-wise~\cite{saeb2017leakage}. Permutation analyses confirm that subject identity alone can carry the reported performance~\cite{chaibubneto2017identity}. In natural language inference (judging whether one sentence logically follows from another), models score far above chance while shown only half of each input, because the annotators' writing habits leak the label~\cite{gururangan2018annotation,poliak2018hypothesis}. In each case an auxiliary attribute (the patient, the annotator's habits) is genuinely predictive within the benchmark, so the model generalises to the held-out set while possessing none of the claimed capability.

Corporate disclosure has the same structure. Firms file repeatedly across a long panel, and volatility is strongly firm-persistent: a calm utility remains calm, a turbulent young technology firm remains turbulent. An expressive model can therefore earn genuine out-of-sample accuracy by recognising \emph{which firm is speaking} rather than reading \emph{what is being said}, and the chronological split that is the field's standard hygiene does nothing about it, because it cuts dates rather than filers (Figure~\ref{fig:bg-shortcut-routes}). Yu et al.'s same-company evidence (Section~\ref{sec:bg-text}) shows the concern is empirical, not hypothetical~\cite{yu2025samecompany}. The same-entity shortcut moreover has a property most shortcuts lack: it is not spurious. Firm identity really does predict future volatility; the difficulty is that the prediction comes free from the firm's own history, so text deserves no incremental credit for re-supplying it. Purging identity from the text is therefore the wrong response; the credit must instead be re-assigned, scoring text against a reference that already knows the firm, in the limit a firm fixed effect. 

\begin{figure}[tbp]
\centering
\includegraphics[width=\textwidth]{figures/C2_identity_shortcut.pdf}
\caption[Two routes to an accurate volatility forecast]{Two routes, one score. A filing reaches the forecast through its content (Route~C) or by identifying its filer --- names, recycled boilerplate, house style --- whose volatility the price history already gives (Route~I). Held-out accuracy scores the two together and never reports the split.}
\label{fig:bg-shortcut-routes}
\end{figure}


Comparisons against so strict a reference must themselves be sound, and forecast evaluation has accumulated a hygiene guarding claims from two directions. Against false positives, three instruments recur. Placebo tests rerun a pipeline with the claimed signal destroyed, permuting which document is attached to which outcome; any credit still awarded exposes machinery rather than information~\cite{ojala2010permutation}. Multiplicity control answers a cruder hazard. A grid of many model--horizon cells hands roughly one in twenty a nominally significant verdict by chance. Holm's step-down procedure~\cite{holm1979simple} bounds the probability of even one false rejection within a family declared in advance (here, ahead of inspection), and holds under arbitrary dependence. Clustered inference recognises that forecast errors are not independent draws: filings sharing a trading day share one market-wide shock, so the day, not the filing, is the independent unit~\cite{cameron2011twoway}. A fourth guards against false negatives: when a challenger nests the benchmark, noise in its extra parameters biases naive loss comparisons against it, which Clark and West's adjustment corrects~\cite{clarkwest2007}.

Hygiene runs equally against over-reading null results. A null is only as informative as the test's sensitivity, so statistical power (the probability that a true effect of a given size would have been detected) must be stated rather than presumed~\cite{cohen1988power}. The demanding form reports a \emph{minimum detectable effect} beside every null (the smallest true improvement the test would have flagged at conventional power, given the panel's actual noise and dependence). The strongest form calibrates sensitivity empirically, injecting a synthetic signal of known size and measuring how often the apparatus recovers it. Together, shortcut learning and evaluation hygiene define the standard against which the positive results of Section~\ref{sec:bg-text} must be read, and no study there meets it jointly. None combines an identity-aware reference, day-level clustered inference, family-wise multiplicity control, placebo gating and a stated minimum detectable effect in a single design.

\section{Methods and Techniques}
\label{sec:bg-methods}

No surveyed design assembles those instruments, so they are assembled here. This section sets out, neutrally, the option space an incremental-value study must draw on; Section~\ref{sec:bg-choice} then commits to choices.

\subsection{The HAR Family}
\label{sec:bg-methods-har}

Corsi's HAR-RV~\cite{corsi2009har} (Section~\ref{sec:bg-price}) anchors the price side, but it is the floor of a family rather than its ceiling. SHAR~\cite{pattonsheppard2015shar} splits the daily regressor into the signed semivariances introduced there, HARQ~\cite{bollerslev2016harq} shrinks the weight on the daily lag on days when it is imprecisely measured, and HAR-X variants append exogenous predictors, most naturally an option-implied volatility index~\cite{blair2001forecasting}.

\subsection{Losses under a Noisy Proxy}
\label{sec:bg-methods-loss}

Volatility is latent, so forecasts are scored against realised volatility, itself a noisy estimate (a \emph{proxy}) of the true quantity. Patton~\cite{patton2011proxies} showed that many intuitive loss functions can reverse forecast rankings when a proxy stands in for the truth, and that only a restricted family, containing squared error and the quasi-likelihood loss QLIKE, ranks forecasts consistently regardless of proxy noise. QLIKE, $y/f - \ln(y/f) - 1$ for proxy $y$ and forecast $f$, penalises relative rather than absolute error and is markedly harsher on under-prediction, the economically dangerous direction for risk management. One subtlety is routinely left implicit: QLIKE may take volatility or variance as its argument. Both conventions are proxy-robust but sit on different numeric scales, so levels must never be compared across conventions.

\subsection{Inference: Dependence, Clustering, Multiplicity}
\label{sec:bg-methods-inference}

The canonical instrument for comparing two forecast sequences is the Diebold--Mariano (DM) test~\cite{diebold1995comparing}. The test forms the loss differential between the forecasts and asks whether its mean differs from zero, with a heteroskedasticity- and autocorrelation-consistent (HAC) variance~\cite{neweywest1987} absorbing the serial dependence that overlapping multi-day targets induce, and a small-sample correction at longer horizons~\cite{harvey1997testing}. In a panel of filings the day-level dependence of Section~\ref{sec:bg-shortcut} dominates, and a test counting every filing as independent evidence overstates precision. The panel-econometric remedies are to cluster the differential by trading day (average within days, then test the daily series), or two ways, by firm and day simultaneously, the conservative extension of Cameron, Gelbach and Miller~\cite{cameron2011twoway}. The Clark--West adjustment~\cite{clarkwest2007} (Section~\ref{sec:bg-shortcut}) applies when a challenger nests the reference. Multiplicity across model--horizon cells is handled by Holm's step-down procedure~\cite{holm1979simple}, which controls the family-wise error rate without independence assumptions and with uniformly more power than the Bonferroni correction it refines. False-discovery-rate control~\cite{benjamini1995controlling} is the standard alternative. Analytic tests are complemented by the block bootstrap~\cite{politis1994stationary}, which resamples contiguous blocks so that dependence is preserved within each draw. Under day-clustering the natural block is the trading day. Set-level procedures ask joint questions: whether any model in a set beats a benchmark (the test for superior predictive ability~\cite{hansen2005spa}) and which subset remains statistically indistinguishable from the best (the model confidence set~\cite{hansen2011mcs}).

\subsection{Encompassing, Combination, Recalibration}
\label{sec:bg-methods-combination}

Those tests rank forecasts. A separate tradition asks whether one adds anything to another, and answers in three ways. Encompassing tests~\cite{harvey1998encompassing} ask whether one forecast's error is reduced by the other. A significant weight on the second means the first does not \emph{encompass} it; the regression detects extra information without locating its source. Bates and Granger~\cite{batesgranger1969combination} went further, showing that a weighted combination frequently outperforms either component alone. A forecast that loses badly on its own can still lower a combination's loss, provided its errors are not the reference's, so a standalone verdict and an incremental one need not agree. Mincer--Zarnowitz regressions~\cite{mincer1969evaluation} supply the third strand, projecting the outcome on the forecast alone; unit slope and zero intercept mark a forecast trustworthy as issued, and the fitted coefficients double as a recalibration. The three strands assemble into one incremental-value instrument. Nest reference and candidate in a small regression, fitted on a held-out period and frozen thereafter, so the reference is recalibrated before the candidate is credited and its coefficient isolates the marginal contribution. Fitting in logarithms suits a positive, right-skewed target. Combination weights estimated on any part of the test period turn evaluation into training.

\subsection{Fine-Tuning, Frozen Embeddings and Prompting}
\label{sec:bg-methods-llm}

What those instruments compare still has to be produced. A pretrained Transformer can forecast in three regimes, ordered by how much of it the task moves. \emph{Fine-tuning} (C1--C4 in Table~\ref{tab:model-matrix})~\cite{devlin2019bert} moves every weight, buying maximal adaptation at maximal risk. Against a noisy scalar target the capacity-to-signal ratio is extreme, and the random seed (the integer fixing weight initialisation and data ordering) alone makes a single run one draw from a wide distribution~\cite{dodge2020finetuning}. Averaging across independently trained replicas, a \emph{seed ensemble}, is the control. \emph{Frozen embeddings} (C5) leave the encoder untouched, a feature extractor beneath a light regularised head, trading adaptation for stability and admitting models far too large to fine-tune. \emph{Prompting} (C6, Qwen3-32B) trains nothing, asking an instruction-tuned model in natural language for a numeric forecast~\cite{lopezlira2023chatgpt}. Its economy carries the three threats of Section~\ref{sec:bg-prompted}, contamination, elicitation fragility and instrument degeneracy, each demanding dedicated probes.

\section{Choice of Methods}
\label{sec:bg-choice}

The choices below are made to keep three claims apart: that an increment is detectable, that it is attributable to content rather than to the filer, and that it is realisable as economic value.

\textbf{The reference, and the ladder above it.} HAR-RV~\cite{corsi2009har} is adopted as the primary reference: the canonical parsimonious model in the realised-volatility tradition, and one that shares a least-squares protocol with every text challenger, so comparisons isolate the information source rather than the estimation machinery. The weak-baseline critique of Section~\ref{sec:bg-survey}, however, makes a single \emph{raw} reference untenable, for two reasons the project's own data confirm. First, raw HAR forecasts under-respond on this panel: the label regressed on them in logs has a mean slope of 1.10 % src: results/tables/recal_slope_logspace.md (mean b_log=1.102;
% forecast_combination.md FAMILY 3's recal_b=1.44 is the level-space slope)
where a calibrated forecast would earn one, and an uncorrected reference donates its miscalibration to any challenger as spurious credit. Second, the HAR extensions surveyed above imply that a gain over plain HAR-RV may be price-side information in disguise. Text is therefore audited against a \emph{reference ladder}: the recalibrated HAR as the most generous rung; a firm-identity-augmented reference that additionally absorbs each firm's own average volatility; and a maximal pool combining five price models % src: writing/dissertation/supporting/sections/05_protocol.tex (reference interval: HAR, SHAR, GARCH, EGARCH, ARIMA)
(HAR, SHAR, GARCH, EGARCH, ARIMA), with an equal-weight variant guarding against overfitted pool weights. Where an increment stops on this ladder is what separates a detectable gain from an attributable one.

\textbf{The model spectrum.} A null's meaning depends on the spectrum that produced it. Had only a prompted large language model been tested, a null could not separate text that carries nothing from a model that finds nothing. Had only bag-of-words, the verdict would concern bag-of-words. The breadth required by Objective~O2, from word counts to a prompted 32-billion-parameter model, is therefore not a survey but the condition under which a negative result can be read. % src: frozen 04_methods.tex Table tab:models (B1--B4 word counts and dictionaries; C1--C4 fine-tuned encoders; C5 frozen 7--8B embedders; C6 prompted Qwen3-32B)
Five long-document strategies enter because long-form filings exceed every fine-tuned window (Section~\ref{sec:bg-neural}). A byte-identical parity arm separates prompting from excerpt curation, and cross-family probes ask whether an increment survives outside one lineage. % src: frozen 03_dataset.tex Table 1 (100% of long-form filings above 4,096 tokens); frozen 04_methods.tex (S1--S5; C5x input-parity arm); frozen 07_ablations.tex (cross-family probes)
The bound is plain: the null covers the classes the spectrum contains, not an architecture nobody ran.

\textbf{The loss.} QLIKE is primary for the combination grid and the standalone leaderboard, per horizon and per disclosure channel, because it is robust to proxy noise and penalises exactly the under-prediction of risk that matters economically~\cite{patton2011proxies}; the 180-comparison census is scored on squared error, its QLIKE restatements reported alongside. The unit convention is declared rather than left implicit: combination cells are scored in volatility units, the standalone leaderboard in variance units, and levels are never compared across the two conventions. % src: writing/dissertation/supporting/sections/05_protocol.tex (counting convention)

\textbf{The incremental instrument.} The central test is a validation-fitted log-space forecast combination in the Mincer--Zarnowitz and Bates--Granger tradition~\cite{batesgranger1969combination,mincer1969evaluation}. The reference and text forecasts are nested in a single exponential-affine regression whose weights are estimated on the validation years alone and frozen before the test period (Chapter~\ref{ch:data} states the equation). Prediction-level combination poses the incremental question most directly: standalone accuracy cannot distinguish text that carries information from text that merely inherits volatility persistence. The combiner credits text only for what it adds beyond a recalibrated reference.

\textbf{Identity controls.} Firms file repeatedly, which makes the same-entity shortcut of Section~\ref{sec:bg-shortcut} this panel's central confound. Three controls separate content from identity: anonymisation (re-scoring filings with identifying language removed), a matched-firm swap (each filing replaced by a comparable firm's), and a zero-text reference augmenting recalibrated HAR with each firm's mean volatility. Each asks how much of an advantage is attributable to the \emph{what} once the \emph{who} is removed.

\textbf{The inference basis.} Two choices were promoted to primary during the investigation, on measured evidence rather than convention. Every trained text and fusion model ran under three seeds (2026--2028), % src: writing/dissertation/supporting/sections/05_protocol.tex (seeds; prompted arms have no training seed to vary)
and single-seed results proved fragile: 37 of 144 trained text and fusion cells flip the sign of their DM statistic across seeds. % src: writing/dissertation/supporting/sections/05_protocol.tex (seed fragility)
The per-observation seed-ensemble forecast is therefore the primary basis. Filings cluster in calendar time (10 to 25 per trading day share one market shock), % src: writing/dissertation/supporting/sections/05_protocol.tex (inference paragraph)
and day clustering shrinks each $|\mathrm{DM}|$ to a median 0.31 of its observation-level value over the 180 standalone comparisons, and to 0.49 over the single-seed 69-cell combination grid. % src: results/tables/dm_pairwise_clustered.md (0.31); results/tables/m1_clustered.md (0.49)
The test of record is therefore the day-clustered DM test on the daily differential series (794--996 trading days per panel), % src: writing/dissertation/supporting/sections/05_protocol.tex; results/tables/dm_pairwise_clustered.md (n_days)
with Holm correction within pre-declared families~\cite{holm1979simple}. Two further gates discipline the counts (Section~\ref{sec:bg-shortcut}). At the primary rung, a label-shuffle placebo replaces the text forecast with a permuted copy, so that any credit it still earns is machinery rather than information. A power calibration injects a synthetic signal of known size and reports the minimum detectable effect beside every null. Chapter~\ref{ch:data} specifies the full protocol; Chapter~\ref{ch:validation} evidences its behaviour under stress.
